\section{Síntesis y ampliación}

\subsection{Un hilo conductor: la existencia del conocimiento no equivale a su disponibilidad}

El concepto más importante de esta clase no es un solo benchmark de reversal, sino dividir el "aprendizaje" en cuatro capas: la experiencia proporciona contenido explícito; el contenido contiene información latente; el proceso de formación genera representaciones internas; y durante la prueba, la interfaz determina cuánta de ella se convierte en información utilizable. La diferencia entre parámetros y contexto ocurre precisamente en las dos últimas capas.

\begin{importantbox}{Marco final del conocimiento}
\begin{enumerate}
  \item El aprendizaje paramétrico y el ICL pueden absorber el mismo hecho, pero generan comportamientos de acceso y generalización distintos.
  \item Reversal, silogismo y codebook son sondas controladas para probar la relación latente; no son agentes completos de toda inteligencia.
  \item La coocurrencia puede mejorar el rendimiento real, pero también ocultar la falta de inferencia sistemática.
  \item El aumento en contexto (in-context augmentation) convierte las implicaciones latentes en supervisión explícita de forma offline.
  \item La recuperación episódica restaura la experiencia original al contexto cuando aparece el objetivo, aunque la recuperación real sigue siendo un problema difícil.
  \item El aprendizaje por regeneración RL reconstruye información en tiempo de prueba, permitiendo transferencia de estructura parcial, pero sigue estando limitado por la búsqueda ante true reversal.
  \item Las tres rutas deben combinarse según costo de entrenamiento, costo de inferencia, cobertura, trazabilidad y recuperación ante fallos.
  \item La analogía con la inteligencia natural respalda la necesidad computacional de sistemas de memoria complementarios, pero no demuestra una implementación neuronal idéntica.
\end{enumerate}
\end{importantbox}

\subsection{De los resultados del aula a la arquitectura del sistema}

Un sistema práctico puede adoptar una estrategia estratificada: realizar aumento offline para relaciones latentes de alta frecuencia, estables y verificables; usar recuperación trazable para hechos de cola larga y experiencias de usuario; emplear razonamiento limitado en tiempo de prueba para tareas que no pueden recuperarse directamente pero tienen estructura general; y activar aclaración, llamadas a herramientas o escalado humano cuando la confianza en la recuperación es baja o la búsqueda de inferencia es excesivamente larga.

\begin{lstlisting}[language=Python,caption={Estrategia de decisión para desplegar tres puentes combinados}]
def answer(query):
    direct = parametric_model.answer(query)
    if verifier.confident(direct):
        return direct

    memories = retriever.search(query)
    if retriever.has_high_recall_evidence(memories):
        return grounded_model.answer(query, context=memories)

    reasoning = reasoning_policy.solve(query, token_budget=budget(query))
    if verifier.confident(reasoning.answer):
        return reasoning.answer

    return request_clarification_or_human_review(query)
\end{lstlisting}

\begin{knowledgebox}{Matriz de evaluación recomendada}
No se debe medir solo QA explícita. Se deben realizar pruebas cruzadas al menos en forward, reversal, multi-hop, alternative goal, cross-lingual, long-context distractor, retrieval miss, conflicting memory y reasoning limitado por cómputo, reportando por separado accuracy, latencia, tokens, recall, completitud de procedencia y calibración.
\end{knowledgebox}

\subsection{Preguntas de investigación aún abiertas}

Primero, ¿cómo entrenar representaciones paramétricas para que soporten naturalmente el indexado en múltiples direcciones sin depender de la enumeración? Segundo, ¿es posible enseñar al modelo a decidir "qué escribir en parámetros, qué guardar como episodio o qué solo requiere contexto a corto plazo"? Tercero, ¿cómo entrenar un retriever general para que el recall no dependa de plantillas de tareas conocidas? Cuarto, ¿puede el razonamiento RL aprender operaciones internas basadas en recuperación, al mismo tiempo que limita los tokens y la confianza errónea? Quinto, ¿cómo convertir las fugas de prompt, alucinaciones y deriva del aumento en métricas auditables?

\begin{warningbox}{Tres conclusiones que no se pueden extraer de esta clase}
No se puede concluir que el ICL siempre es superior al fine-tuning; no se puede concluir que todo el conocimiento paramétrico sea incapaz de realizar razonamiento relacional; no se puede concluir que el hipocampo sea equivalente al RAG. La clase demuestra diferencias en la disponibilidad de información bajo condiciones controladas específicas y propone tres métodos de puente verificables y combinables.
\end{warningbox}

\subsection{Lecturas adicionales}

Se recomienda leer siguiendo la cadena de evidencia: primero, leer el Reversal Curse de Berglund et al., para comprender el fallo en el acceso direccional; luego, estudiar el estudio controlado sobre ICL/fine-tuning de Lampinen et al., verificando controles de contexto en todo el conjunto de datos, silogismos y desde cero (from-scratch); después, leer Latent Learning, prestando especial atención al oráculo episódico y la memoria complementaria; finalmente, leer Improving Latent Generalization Using Test-time Compute para revisar la transferencia RL y las limitaciones de reversal. Sobre la relación entre ICL y optimización, se pueden contrastar los resultados teóricos limitados de Akyürek et al. con los de Dai et al.

\begin{knowledgebox}{Snapshot de la clase}
Esta clase tuvo lugar el 2026-05-07. Los apuntes utilizan las versiones arXiv ya públicas en ese momento: controlled study v3, Latent Learning v3, test-time compute v1 y Reversal Curse v4. Las revisiones posteriores pueden servir como nueva evidencia, pero no deben extrapolarse retroactivamente como resultados que se mostraron el día de la clase.
\end{knowledgebox}

En última instancia, esta clase reformula "cuánto recuerda el modelo" en una pregunta de sistema más precisa: \textbf{¿Puede el modelo reconstruir las experiencias pasadas como información computable actual cuando cambian los objetivos?} Los parámetros, el contexto, la memoria episódica, el aumento y el RL no son cinco módulos aislados, sino divisiones de trabajo diferentes para el mismo ciclo de vida de la información.

\end{document}
